\section{Related Work}
\label{sec:related}

\subsection{Task Formulations in Dental Image Analysis}
\label{subsec:tasks}

\textbf{Single-Task vs Multi-Task Learning.} Most existing dental imaging systems address single tasks in isolation---either tooth segmentation~\cite{li2021deep}, landmark detection~\cite{ali2025interpretable}, or disease classification~\cite{lee2017qlf,zhang2020dualchannel}. While these approaches achieve strong performance on their specific objectives, they require deploying multiple separate models for comprehensive clinical assessment, increasing computational overhead and integration complexity. Recent work has begun exploring multi-task frameworks: Ali et al.~\cite{ali2025interpretable} combined landmark detection with malocclusion classification using Vision Transformers, demonstrating that shared representations can improve generalization. However, \textit{lightweight hybrid architectures for multi-class oral disease classification remain underexplored}, particularly designs that effectively balance local texture extraction with global spatial reasoning while maintaining deployment feasibility on standard clinical workstations.

\textbf{Single-Label vs Multi-Label Classification.} Early CNN-based methods applied single-label classification for caries or plaque detection on QLF and dual-channel dental images~\cite{lee2017qlf,zhang2020dualchannel}. Segmentation-to-classification pipelines have been proposed to improve localization~\cite{nguyen2022caries}. However, single-label approaches cannot handle multiple concurrent diseases, which is common in clinical practice. Multi-label classification addresses this limitation by predicting multiple labels simultaneously, as demonstrated on DENTEX~\cite{hamamci2023dentex} and pediatric/adult panoramic datasets~\cite{alsakar2024multilabel}. While multi-label formulations better reflect clinical reality, they introduce challenges in model calibration, class weighting, and annotation cost~\cite{lof2ml2025}. In this work, we adopt a multi-class formulation since the Oral Diseases dataset provides single-label annotations where each image is assigned exactly one disease category.

\subsection{Architectural Approaches}
\label{subsec:architectures}

\textbf{CNN-Based Methods.} Convolutional neural networks have been the dominant paradigm for dental image analysis. U-Net and its variants are widely adopted for tooth segmentation~\cite{hassanein2024utilizing}, leveraging encoder-decoder architectures with skip connections to preserve spatial information. ResNet and DenseNet backbones provide strong feature extraction for classification tasks~\cite{george2025deep}. While CNNs excel at capturing local structures (tooth boundaries, bone edges), their limited receptive fields struggle with long-range dependencies required for global anatomical context~\cite{li2021deep}. This limitation is particularly problematic for landmark detection tasks where spatial relationships between distant points (e.g., bilateral symmetry, maxillary-mandibular correlations) are critical.

\textbf{Transformer-Based Methods.} Vision Transformers have recently emerged in dental imaging, with Ali et al.~\cite{ali2025interpretable} demonstrating that ViT-based architectures achieve superior performance on cephalometric landmark detection compared to CNNs, attributed to their global attention mechanisms. However, pure Transformer models require large training datasets and exhibit higher computational complexity ($\mathcal{O}(N^2)$ attention), limiting their applicability in data-scarce clinical settings~\cite{wu2025abductive}.

\textbf{Hybrid Architectures.} Recognizing the complementary strengths of CNNs and Transformers, recent work has explored hybrid designs. TransUNet and SegFormer combine CNN encoders with Transformer decoders for medical image segmentation~\cite{araujo2025twostep}. Swin Transformer introduces hierarchical shifted window attention for linear complexity~\cite{george2025deep}. While promising, existing hybrid methods for dental imaging have not been optimized for multi-task learning or lightweight deployment.

\textbf{Two-Step Pipelines.} Multi-stage architectures separate detection/ROI extraction from segmentation/classification to improve accuracy and handle class imbalance. Lin et al.~\cite{lin2023mandibular} developed a coarse-to-fine 3D-UNet pipeline for mandibular canal segmentation from CBCT images, where coarse localization reduces search space before fine segmentation. Kim et al.~\cite{kim2025multistage} combined YOLOv5 tooth detection with Attention U-Net for caries segmentation and staging on panoramic X-rays. While two-step pipelines achieve strong accuracy, they increase deployment complexity by requiring multiple models and suffer from error propagation when the first stage fails.

\subsection{Optimization Techniques for Clinical Deployment}
\label{subsec:optimization}

\textbf{Class Imbalance Handling.} Dental datasets exhibit severe class imbalance, with rare tooth types (e.g., wisdom teeth) and pathologies significantly underrepresented. Without proper handling, models bias toward majority classes and miss clinically important rare conditions~\cite{hamamci2023dentex,sharma2024oralhealth}. Common mitigation strategies include: (1) \textit{Data-level techniques}---oversampling (SMOTE), undersampling, or targeted augmentation for minority classes~\cite{hassanein2024utilizing}; (2) \textit{Algorithm-level techniques}---class weighting, focal loss, or hybrid loss functions (LMFLOSS) that emphasize hard samples~\cite{li2022imbalanced}; (3) \textit{Cascade detection}---two-step pipelines that first detect abnormal regions before classification, reducing imbalance effects~\cite{jiang2021imbalanced,araujo2025twostep}. While these techniques improve minority class sensitivity, they often increase training complexity and require careful hyperparameter tuning.

\textbf{Lightweight Model Design.} Clinical deployment requires models that can run on standard workstations without high-end GPUs. Gholami et al.~\cite{gholami2025efficient} survey compression techniques including pruning, quantization, and knowledge distillation. Canali et al.~\cite{canali2022challenges} demonstrate that efficient attention mechanisms (e.g., Bi-level Routing Attention) can reduce computational complexity while maintaining performance. However, most existing dental imaging models prioritize accuracy over efficiency, with limited exploration of systematic lightweight optimization. This gap hinders practical adoption in resource-constrained healthcare facilities.

\subsection{Evaluation and Robustness in Real-World Settings}
\label{subsec:evaluation}

Ensuring model stability under real-world variations---different imaging equipment, lighting conditions, patient populations, and noise levels---is critical for clinical adoption. Studies on intraoral images using segmentation-to-classification pipelines have tested models under various imaging conditions, combining augmentation to improve noise resistance~\cite{nguyen2022caries}. Multi-label datasets such as DENTEX~\cite{hamamci2023dentex} and pediatric/adult panoramic datasets~\cite{alsakar2024multilabel} have been used to evaluate generalization across diverse patient populations and concurrent diseases. However, most evaluations remain limited to single-institution datasets with consistent imaging protocols, providing limited evidence of cross-domain generalization. Speight et al.~\cite{speight2017screening} highlight the need for multi-center validation to assess model robustness across different clinical settings, yet such comprehensive evaluations remain rare in dental imaging literature.

\subsection{Comparison of Representative Methods}
\label{subsec:comparison}

Table~\ref{tab:comparison} summarizes key characteristics of representative dental imaging methods. Existing approaches predominantly focus on single tasks, with limited exploration of multi-task frameworks that jointly address segmentation and landmark detection. While hybrid architectures show promise, most implementations prioritize accuracy over computational efficiency, hindering deployment in resource-constrained clinical environments. Furthermore, evaluation protocols vary significantly across studies, making direct performance comparison difficult.

\begin{table*}[t]
\centering
\caption{Comparison of representative methods for oral disease classification. Our approach achieves competitive performance with significantly fewer parameters than pure Transformer architectures while incorporating global attention mechanisms absent in pure CNN approaches.}
\label{tab:comparison}
\begin{tabular}{lccccc}
\hline
\textbf{Method} & \textbf{Task} & \textbf{Architecture} & \textbf{Params} & \textbf{FLOPs (G)} & \textbf{Limitations} \\
\hline
ResNet-50~\cite{he2016deep} & Classification & CNN & 25.6M & 4.1 & Limited global context \\
EfficientNet-B3~\cite{tan2019efficientnet} & Classification & CNN & 12.0M & 1.8 & No explicit attention mechanism \\
ViT-Base~\cite{dosovitskiy2021vit} & Classification & Transformer & 86.0M & 17.6 & High computation; data-hungry \\
Swin-Tiny~\cite{liu2021swin} & Classification & Transformer & 28.3M & 4.5 & Still computationally expensive \\
George et al.~\cite{george2025deep} & Multi-label & ResNet-50 & 25.6M & 4.1 & Pure CNN; no global attention \\
\hline
\textbf{Ours (Base)} & \textbf{Classification} & \textbf{Hybrid CNN-ViT} & \textbf{15.7M} & \textbf{2.8} & \textbf{Balanced local-global features} \\
\hline
\end{tabular}
\end{table*}


\subsection{Research Gap and Positioning}
\label{subsec:gap}

Based on the comprehensive literature review, we identify three critical gaps in existing dental image analysis research:

\textbf{Gap 1: Insufficient Integration of Local-Global Features.} Existing approaches for oral disease classification rely primarily on either pure CNNs (limited global context) or pure Transformers (high computational cost, data-hungry). Hybrid architectures for oral disease classification remain underexplored, particularly lightweight designs that balance local texture extraction with global spatial reasoning while maintaining deployment feasibility on standard clinical workstations. This gap limits the ability to capture both fine-grained lesion boundaries and long-range spatial dependencies crucial for accurate disease recognition.

\textbf{Gap 2: Insufficient Focus on Lightweight Deployment.} While recent hybrid architectures achieve impressive accuracy, they typically require high-end GPUs (e.g., NVIDIA V100, A100) for inference, limiting adoption in resource-constrained healthcare facilities. Pure Transformer models like ViT-Base (86M parameters) and even efficient variants like Swin-Tiny (28.3M) remain computationally expensive for real-time clinical deployment. This gap hinders practical adoption, particularly in underserved regions with limited computational infrastructure.

\textbf{Gap 3: Limited Robustness Evaluation.} Most studies evaluate models on single-institution datasets with consistent imaging protocols, providing limited evidence of generalization to real-world variations in equipment, image quality, and patient populations. Cross-domain validation and robustness analysis under noise, artifacts, and anatomical variability are rarely reported, raising concerns about reliability in diverse clinical settings.

\textbf{Our Contribution.} This work addresses these gaps by proposing a lightweight hybrid Transformer-CNN architecture for multi-class oral disease classification from intraoral images. Our approach integrates: (i) window-based self-attention for global context modeling, (ii) learnable gated cross-attention fusion enabling adaptive CNN-Transformer feature balancing, and (iii) compact design (15.7M parameters) suitable for clinical deployment on standard hardware. Our base model achieves competitive performance while being 82\% smaller than ViT-Base (86M), 45\% smaller than Swin-Tiny (28.3M), and comparable to EfficientNet-B3 (12M) while incorporating global attention mechanisms absent in pure CNN architectures. This enables deployment on standard clinical GPUs (8GB VRAM) with inference time $<$50ms per image. Our evaluation includes extensive ablation studies and robustness analysis under synthetic perturbations, providing comprehensive evidence of clinical viability.
